% AI-written, AI cross-reviewed; no human sign-off; not a claim that RH is proved
% Verbatim sections of local blueprint v13. See README.md for macros, edition and source spans.

% blueprint/ch_P.tex:1984-2417
\subsection{Lattices, signs and the two duals}
Let $K$ be totally complex, $n=[K:\mathbb Q]$, $\mathcal O=\mathcal O_K$,
and $a\mapsto a^*$ an involution preserving $\mathcal O$. For the distinguished
embedding $\sigma_0$ representing $v_0$ assume
$\sigma_0(a^*)=\overline{\sigma_0(a)}$ for every $a\in K$.
At other embeddings the analogous identity is required for entries in the
specified subfield $K_0$. Write $\bar y$ for the conjugate of a finite
place $y$, $v_y$ for the exponent in the principal fractional ideal, and
$w_y=\log\mathrm N y/n$. The involution permutes the primes, preserves their
residue field sizes, and gives $v_y(a^*)=v_{\bar y}(a)$.
In particular $w_y=w_{\bar y}\ge0$. For $x\ne0$, factoring $(x)$ and taking
the field norm in its embeddings gives the normalization
\[
   \sum_y w_y v_y(x)=\frac{1}{n}\log|N_{K/\mathbb Q}x|
                  =\frac{2}{n}\sum_{v\mid\infty}\log|\sigma_v(x)|.
\]
All finite-place sums in this formula have finite support. Complex places
are counted once on the right.
See \path{IsCMSub}, \path{IsConjAt} and \path{lw}
(\path{QuotientMetricDivisors/Statement.lean:18--64}),
\path{ValBasic} (\path{QuotientMetricDivisors/ProdFormula.lean:41--47}) and
\path{prod_formula'} (\path{HermitianTwistDivisors/ProdFormula.lean:39--93}).
Indeed the involution sends the factorization of $(a)$ to that of $(a^*)$
and induces an isomorphism on each corresponding residue field. Applying
it twice fixes each prime, so reindexing the finite support leaves its
sum unchanged. This proves the conjugation rules used below, including
at a fixed prime (\path{finite_sum_conjugation}, \path{ValBasic.lean:147--152,170--173}).

For full lattices $A,B$ in one $K$-space, take the ordinary localization
at $y$ and set
\[
 i_y(A,B)=\operatorname{length}(A_y/(A\cap B)_y)
                    -\operatorname{length}(B_y/(A\cap B)_y).
\]
Using a common sublattice in both terms gives the same difference by
additivity of length. This also proves $i_y(A,C)=i_y(A,B)+i_y(B,C)$.
If a basis for $B_y$ is obtained from one for $A_y$ by right multiplication
by $C$, then $i_y(A,B)=v_y(\det C)$. Indeed, put the matrix into elementary
divisor form over the discrete valuation ring after clearing denominators.
In one coordinate the two lengths have difference $r$ for multiplication by
a uniformizer to exponent $r$, including $r<0$. Summing the exponents gives
the assertion. A semilocal basis at $\{y,\bar y\}$ is an ordinary basis
after localization at either prime, since the denominators inverted in
the semilocal ring are units at both. Applying this calculation to each
localization gives both clauses, including $y=\bar y$. The source proves
the determinant formula by diagonal and transvection reduction. See \path{Formalism.HermitianLatticesTheory.IndexBasic}
(\path{HermitianLattices/StatementLat.lean:192--203}) and
\path{index_change_formula} (\path{IdxChange.lean:640--647,671--694}).

For a Hermitian form $h(x,z)=x^*Hz$, write
$T^\#=\{z:h(T,z)\subset\mathcal O\}$. The index correction is
$j_y(H,T)=i_y(T,T^\#)$. The dual that occurs in the final slope is the
$\mathcal O$-linear evaluation dual in the $K$-dual space, equipped with
the dual frames. These constructions have different roles. Given a frame
$X_v$, its positive metric matrix is $A_v(X)=(X_v^{-1})^*X_v^{-1}$.
For columns of a matrix $f$, its complex Gram is
$\sigma_v(f)^* A_v(X)\sigma_v(f)$, and its rational Gram for $H$ is $f^*Hf$.
We use ordinary determinants, so that the change of determinant is a squared
absolute value at a complex embedding.
See \path{Formalism.DivisionValues.gramR}, \path{detRatio}, \path{jw}
(\path{QuotientMetricDivisors/Statement.lean:38--93}) and
\path{Formalism.HN.DualPl} (\path{HarderNarasimhan/QuotRealize.lean:72--82}).

\begin{lemma}[The signed local correction]\label{lem:P5-dual}
\lean{Formalism.DivisionValues.conjugate_weight_gram_general}\uses{}
If $H^*=H$ is nonsingular, $T$ is full, and the columns $f$ form a basis
of $T$ after semilocalizing at $y,\bar y$, then
\[
                 j_y(H,T)=-v_y\bigl(\det(f^*Hf)\bigr).
\]
\end{lemma}
\begin{proof}
The semilocal ring inverts elements outside both primes and is stable under
the involution. Its Dedekind factorization has at most two nonzero primes,
so every ideal is principal (choose an element with just the specified
orders at these primes by approximation, or by the Chinese remainder
theorem). A finitely generated torsion-free module is free there. A full
lattice therefore admits such a basis, with $K$-linearly independent
columns. Semilocalization of its dual is the integral dual after
semilocalizing: to clear the finitely many pairings against a generating
set one may use a common denominator outside both primes. Consequently the
columns $f(f^*Hf)^{-1}$ span the semilocal dual. Indeed their pairings with
$f$ form the identity, and the pairings of any vector with $f$ are exactly
its coefficients in that dual basis. Ordinary localization at $y$ now
gives the preceding determinant formula with matrix $(f^*Hf)^{-1}$.
Nondegeneracy and independence ensure the inverse exists.
This proves the sign even when the two primes coincide or $T$ fails to be
integral for $h$.
See \path{semilocal_basis_exists} (\path{LocBasis.lean:110--143}),
\path{jwGramW_of} and \path{conjugate_weight_gram_general}
(\path{ConjugateWeightLatticeOfGeneral.lean:18--124}).
\emph{Status: Complete.}
\end{proof}

\subsection{The real lattice and its evaluation dual}

\begin{lemma}[Covolume in a complex frame]\label{lem:P5-gram-realize}
\lean{Formalism.HN.height_gram_determinant}
Let $K$ be totally complex of degree $n=2f$, and let $A$ be a full,
finitely generated $\mathcal O_K$-lattice in the $D$-dimensional space
$V$. In the invertible frames $X_v$, its realization $\rho_X A$ is a
Euclidean lattice of rank $nD$ and positive covolume $C_X(A)$.
If $e_1,\ldots,e_D$ is a $K$-basis in $A$ and $B=\sum_i\mathcal O_K e_i$,
then, writing $\Delta=|\operatorname{disc}K|$,
\[
 [A:B]\operatorname{ht}(X,A)=\operatorname{ht}(X,B)
 =\prod_{v\mid\infty}\det\operatorname{Gram}_v(X,e),
 \qquad \operatorname{ht}(X,A)=2^{fD}\Delta^{-D/2}C_X(A).
\]
\end{lemma}
\begin{proof}
Fix an integral basis $\omega_1,\ldots,\omega_n$ of $\mathcal O_K$.
Take the real matrix $M$ whose $i$th column consists of the real and
imaginary parts of $\sigma_v(\omega_i)$, one embedding per pair.
Replacing each pair of rows $(u,v)$ by $(u+iv,u-iv)$ multiplies the
absolute determinant by 2. The resulting matrix $E$ contains all
embeddings. The $(i,j)$ entry of $E^{t}E$ is
$\operatorname{Tr}_{K/\mathbf Q}(\omega_i\omega_j)$, so
$|\det M|=2^{-f}\sqrt{\Delta}$. This is the unit-ideal case of
Milne \cite[Prop.~4.25]{P-Milne-Cheb} with the real coordinate measure.
See \path{HeightGram.abs_det_amat} (\path{HtGram.lean:218--243}).

Write $U_v$ for the complex matrix whose columns are the coordinates
of $e_j$ in the frame, obtained by applying $X_v^{-1}$ to their
embedded coordinates. Before these matrices act, the columns for
the $\omega_i e_j$ give $D$ copies of $M$, up to row and column
permutations. The real determinant of $U_v$ is $|\det U_v|^2$:
after complexification the real linear map is conjugate to
$\operatorname{diag}(U_v,\overline{U_v})$. The same number is
$\det(U_v^*U_v)>0$. Taking determinants of the composite gives
\[
 C_X(B)=(2^{-f}\sqrt{\Delta})^D
           \prod_{v\mid\infty}\det(U_v^*U_v).
\]
See \path{HeightGram.det_realif}, \path{HeightGram.det_gmat} and
\path{Formalism.HN.htGram_of} (\path{HtGram.lean:85--137,191--197,341--409}).

Each of the finitely many generators of $A$ has coordinates in $K$.
Clearing their denominators with a single positive integer $c$ gives
$B\subset A\subset c^{-1}B$. Thus the middle image is a discrete
subgroup spanning the same real space as $\rho_X B$. In integral
bases for $\rho_X B$ and $\rho_X A$, the absolute determinant of the
inclusion matrix is $[A:B]$ by Smith normal form. It multiplies the
larger lattice's covolume by that index. This proves the formulas,
including $D=0$ with empty determinant one.
See \path{PlaceLattice.exists_denom} and \path{place_lattice_structure}
(\path{DiscreteSpanOfLinIndep.lean:224--248,269--324}) and
\path{height_gram_determinant} (\path{HtGram.lean:411}).
\emph{Status: complete.}
\end{proof}

\begin{lemma}[Height of the evaluation dual]\label{lem:P5-eval-realize}
\lean{Formalism.HN.DualPolygon.ht_dual}\uses{lem:P5-gram-realize}
For $A$ as above, $A^*=\{\phi\in V^*:\phi(A)\subseteq\mathcal O_K\}$
in the dual frame satisfies
$\operatorname{ht}(X^\vee,A^*)\operatorname{ht}(X,A)=1$.
\end{lemma}
\begin{proof}
Write $L=\rho_X A$, $\rho'=\rho_{X^\vee}$, and let $J$ denote
coordinatewise complex conjugation. The trace dual is
$T=\{\phi\in V^*:\operatorname{Tr}(\phi(A))\subseteq\mathbf Z\}$.
The Euclidean dual is exactly $L^\#=2J\rho' T$. Indeed, pairing the
realizations of $x$ and $2J\rho'\phi$ gives
$2\sum_v\operatorname{Re}(\sigma_v\phi(x))=\operatorname{Tr}(\phi(x))$:
the transposed change of coordinates in the dual cancels $X_v^{-1}$
in the primal. Conversely an integral linear form on $L$ extends to
a rational linear form on $V$ because $A$ spans $V$ over $\mathbf Q$.
The map $\phi\mapsto\operatorname{Tr}\circ\phi$ from $V^*$ to the
rational dual is injective (a nonzero $K$-linear form takes the value
$1$, whose trace is $n$) and both rational dimensions are $nD$.
The representing real vectors agree since $L$ spans the real space.
A Euclidean lattice with basis matrix $H$ has dual basis $H^{-t}$.
The factor 2 acts in $nD$ real dimensions and $J$ is an isometry, giving
\[
 \operatorname{covol}(\rho'T)=2^{-nD}C_X(A)^{-1}.
\]
See \path{DualPolygon.ident_trace_eq}, \path{inner_realize_alt} and
\path{dualZ_zl_eq} (\path{DualPl.lean:70--95,114--148,196--312}).

Let $\mathfrak d^{-1}=\{z\in K:\operatorname{Tr}(z\mathcal O_K)
\subseteq\mathbf Z\}$ be the codifferent. Since $\mathcal O_K A=A$,
one has $T=\operatorname{Hom}_{\mathcal O_K}(A,\mathfrak d^{-1})$.
Both modules in $A^*\subseteq T$ are full lattices: in the notation of
the preceding proof $cB^*\subseteq A^*\subseteq B^*$, and $T$ already
has a full real realization. In the integral basis $\omega_i$ the
codifferent has coordinate lattice
$(\operatorname{Tr}(\omega_i\omega_j))^{-1}\mathbf Z^n$.
Hence $|\mathfrak d^{-1}/\mathcal O_K|=\Delta$.
At each nonzero prime $\mathfrak p$, $A_{\mathfrak p}$ is free of rank $D$
over the discrete valuation ring: it is finitely generated and
torsion-free over a principal ideal domain. Localization commutes with
these Hom modules since $A$ is finitely presented. Thus
$(T/A^*)_{\mathfrak p}$ is the direct sum of $D$ copies of
$(\mathfrak d^{-1}/\mathcal O_K)_{\mathfrak p}$. Multiplying residue-field
orders with their composition lengths over the finitely many primes in
the support gives $[T:A^*]=\Delta^D$. This also explains why a global
$\mathcal O_K$-basis of $A$ is unnecessary. See
\path{DualPolygon.dual_inf_ann_eq}, \path{twist_index} and \path{key_lemma}
(\path{DualPl.lean:443--495,689--731,863--872,947--955}).
It follows that
\[
 C_X(A)C_{X^\vee}(A^*)=2^{-nD}\Delta^D .
\]
Each height contributes $2^{fD}\Delta^{-D/2}$, so their product is one
since $2f=n$. This proves the full-lattice clause of
\path{Formalism.HN.DualPolygon.ht_dual} (\path{DualPl.lean:1065--1073}).
\emph{Status: complete.}
\end{proof}

\subsection{The base quotient comparison}
Write $\mathcal V(X,L)=2\log\operatorname{ht}(X,L)/n$, and use
$\langle h_v\rangle=(2/n)\sum_{v\mid\infty}h_v$.
Lemma~\ref{lem:P5-gram-realize} gives the following normalization.
In full rank an independent list $b$ spanning
$B$ over $\mathcal O$ satisfies
$\operatorname{ht}(X,B)=\prod_v\det\operatorname{Gram}_v(X,b)$, and for a
full $L\supseteq B$ one has $\operatorname{ht}(X,L)[L:B]=\operatorname{ht}(X,B)$.
These are \path{Formalism.HN.HtGram} and
\path{height_gram_determinant} (\path{HarderNarasimhan/StatementPl.lean:83--91},
\path{HtGram.lean:411}). The index here is the cardinality of an additive
quotient.

For any full $L$, even if $L$ and $B$ are incomparable, choose an integer
$c\ne0$ with $cB\subset L\cap B$. Comparison of the two indices above $cB$
gives the following covolume formula:
\[
 \mathcal V(X,L)=
 \big\langle\log\det\operatorname{Gram}_v(X,b)\big\rangle
                          +2\sum_y w_y i_y(B,L).                 \qquad\text{(P5.1)}
\]
For a finite torsion module, the logarithm of its size is the sum of the
local lengths times $\log\mathrm N y$: a composition quotient is a single
residue field and contributes only at its own prime. Applying this twice
cancels the common sublattice and proves finite support. Thus changing $L$
to $L'$ in a fixed frame adds $2\sum_yw_y i_y(L,L')$. Applying an
isomorphism $\phi$ to both $L$ and $B$ leaves their local indices unchanged.
If $d_v(\phi)$ denotes its image Gram determinant divided by the source Gram
determinant, the change in $\mathcal V$ between the two frames is exactly
$\langle\log d_v(\phi)\rangle$. The ratio is independent of $b$ because
changing the $K$-basis multiplies both determinants by the same positive
factor. See \path{covolume_via_basis} (\path{LcvBasis.lean:324--427}),
\path{ImageLcv}, \path{IdxLcv} (\path{QuotientMetricDivisors/ProdFormula.lean:54--67}).

\begin{lemma}[The loss at the base type]\label{lem:P5-base}
\lean{Formalism.DivisionValues.reduction_lemma_second}\uses{lem:P5-gram-realize}
Let $P:K^a\twoheadrightarrow K^D$, let $M_*$ be full, and let $H$ have entries
in $K_0$ and be positive definite at every infinite place. Give $Q=P M_*$
the quotient metrics $A_v=(P_v H_v^{-1}P_v^*)^{-1}$ and write
$\Gamma=(PH^{-1}P^*)^{-1}$ over $K$. For semilocal bases $e_y$ of $Q$, set
$G_y=\det(e_y^*\Gamma e_y)$. Then $w_yv_y(G_y)$ has finite support and
\[
                       \mathcal V(X^Q,Q)\le\sum_y w_y v_y(G_y).
\]
\end{lemma}
\begin{proof}
The images of generators of $M_*$ generate $Q$ and their $K$-span
is $P(K^a)=K^D$ since $M_*$ is full and $P$ is surjective
(\path{lemmaRb_of}, \path{ReductionLemmaSecond.lean:163--177}). Write
$\tau(a)=\overline{\sigma(a^*)}$ for the twist of an embedding $\sigma$.
It agrees with $\sigma$ on $K_0$, commutes with conjugation of embeddings,
and twisting twice returns the original embedding. In particular it
permutes the infinite places. If $C=PH^{-1}P^*$, then
\[
          \sigma(C)=P_\sigma H_\sigma^{-1}P_\tau^*.
\]
At $v_0$, both embeddings on the right agree. The right side there is
positive definite, so $\det C\ne0$ in $K$. Away from $v_0$, the two
matrices of $P$ can differ. This is precisely where equality of the rational
and archimedean quotient determinants could fail.
See \path{Rb.tw_tw}, \path{Rb.plc_plc}, \path{Rb.Cmap}, \path{Rb.Hmap_tw}
and \path{lemmaRb_of} (\path{ReductionLemmaSecond.lean:22--41,67--79,185--202}).

For a positive definite complex matrix $S$ and rectangular matrices $U,V$,
put $A=USU^*$, $B=USV^*$, $C'=VSV^*$. If $A$ is invertible, put
$W=V-B^*A^{-1}U$. Direct multiplication gives
$C'=B^*A^{-1}B+WSW^*$. A positive semidefinite matrix increases in
determinant upon adding another one: for a positive definite first matrix
conjugate it to the identity and use the eigenvalues of the second matrix,
then pass to a limit for a singular first matrix. Hence
$|\det B|^2\le\det A\det C'$. If $A$ is singular, a nonzero left null vector
for $A$ also annihilates $U$ by positivity of $S$, hence annihilates $B$.
The same inequality holds in that case. This is
\path{determinant_cauchy_schwarz} (\path{DeterminantCauchySchwarz.lean:67--139}).

Apply this inequality with $U=P_\sigma$, $V=P_\tau$, $S=H_\sigma^{-1}$;
$H_\tau=H_\sigma$ since its entries are in $K_0$. The two real diagonal
Gram determinants are the reciprocals of the metric determinants at the
places of $\sigma,\tau$. Taking reciprocals of positive numbers gives
\[
  \det A_{v(\sigma)}\det A_{v(\tau)}
        \le |\sigma(\det C)|^{-2}=|\sigma(\det\Gamma)|^2 .
\]
Conjugating the representative of a place leaves each positive real
determinant unchanged. Sum logarithms over places: each term on the left
then occurs twice, including the fixed places of the involution.
The product formula gives
$\langle\log\det A_v\rangle\le\sum_y w_yv_y(\det\Gamma)$.
Finally, writing $B=\mathcal O^D$, the two basis changes at $y,\bar y$ give
\[
 v_y(G_y)=v_y(\det\Gamma)+ i_y(B,Q)+i_{\bar y}(B,Q).
\]
Each of these terms has finite support. In the weighted sum the last two
terms combine to $2\sum_yw_y i_y(B,Q)$ by conjugation. Add them to the
archimedean inequality and use (P5.1) in the standard basis.
See \path{Rb.place_ineq}, \path{lemmaRb_of}
(\path{ReductionLemmaSecond.lean:133--157,204--290}).
\emph{Status: complete.}
\end{proof}

\subsection{One type at a time}
Let $t_0$ denote the base type. The types $t_y$ assigned to the primes
satisfy $t_{\bar y}=t_y$, and differ from $t_0$ only at finitely many
primes. For each type we have a full quotient lattice $Q_t$ with frame
$X^Q_t$ and an isomorphism $\iota_t$ to the fixed framed space $E$.
There are finitely many columns $j$, with full $T_{jt}$, forms $H_{jt}$,
frames $X^T_{jt}$ and isomorphisms $\theta_{jt}$ to fixed framed spaces
$W_j$ containing the reference lattices $M_j$. Set
$A_t=\iota_t Q_t$, $B_{jt}=\theta_{jt}T_{jt}$. At $v_0$ each $H_{jt}$
realizes the indicated column metric. It is consequently Hermitian and
nonsingular over $K$, as can be seen after applying the injective embedding
$\sigma_0$. At the base type the column forms have entries in $K_0$ and
realize the metrics at every infinite place. Given semilocal bases $e_{t,y}$
of $Q_t$, we require a $G_{t,y}\in K$ with
$\sigma_0 G_{t,y}=\det\operatorname{Gram}_{v_0}(X^Q_t,e_{t,y})$.
At the base it is the rational Gram of Lemma~\ref{lem:P5-base}.
For the actual quotients one obtains the required number at every type
from the Hermitian $\Gamma_t=(P_tH_t^{-1}P_t^*)^{-1}$, since $\sigma_0$
commutes with the adjoint. Moreover $\det(e^*\Gamma_t e)$ is star-invariant.
These statements are \path{IsHermAt} (\path{RatHermLcvTwist.lean:11--15}),
\path{gram_determinant_quotient_metric}
(\path{ConjTransposeEqOfPosDef.lean:61--90}) and \path{IndefRows.stepB}
(\path{IndefUse.lean:236--269}).

\begin{lemma}[The two required determinant identities]\label{lem:P5-ratio}
\lean{Formalism.L2.Indef.IndefRows.detId}\uses{}
In the realization by indefinite rows the ratios satisfy
\[
  d_v(\iota_t)=\kappa_v^D\prod_j d_v(\theta_{jt})
                  \qquad\text{if }t=t_0\text{ or }v=v_0.
\]
All column dimensions and quotient dimensions equal $D>0$ under the
base-side well-formedness hypotheses.
\end{lemma}
\begin{proof}
At the base one has independent summands spanning the quotient and each
column; the row and column isomorphisms transport them to $E$ and $W_j$.
The parts there have matching dimensions for each label $p$. Pulling them
back along the isomorphisms for any $t$ preserves independence,
exhaustion and dimensions. If their dimensions are $d_p$, summing gives
$\sum_p d_p=\dim E=\dim W_j$ for each $j$. Bijectivity also identifies
$\dim Q_{t_0}$ with $D$ and positivity of the base quotient yields $D>0$.
This proves each rank assertion before any division by $D$.
See \path{IndefRows.part_dim}, \path{IndefRows.dW_eq},
\path{IndefRows.D_pos} and \path{IndefRows.Zpart_decomp}
(\path{IndefUse.lean:109--170}).

Let $N_1$ denote the matrix of a pulled-back target metric and $N_2$ its
source metric. The scalar hypothesis on part $p$ is the bilinear identity
$x^*N_1z=c_p x^*N_2z$ for basis vectors $x$ in this part and every
rational vector $z$, after embedding at $v$. With the concatenated bases
as columns of $C$, it gives the entire matrix identity
\[
 C^*N_1 C=\mathrm{diag}(c_{p(i)})C^*N_2 C .
\]
Consequently $d_v=\prod_p c_p^{d_p}$: the common factor
$|\det C|^2$ cancels. This uses no orthogonality assumption between parts.
For the row take $c_p=\kappa_v\prod_j\lambda_{jtpv}$, and in column $j$
take $c_p=\lambda_{jtpv}$. Interchange the two finite products and use
$\sum_p d_p=D$. The positivity of the frames makes every ratio for a
bijective map strictly positive. The supplied scalar hypotheses hold
in exactly the two cases of the statement. Zero-dimensional summands
contribute exponent zero, and an empty family of columns contributes
the product one.
See \path{Formalism.DivisionValues.det_of_scalar}, \path{spectral_determinant_formula}
(\path{SpecDet.lean:87--132,160--177}) and
\path{IndefRows.detId} (\path{IndefUse.lean:196--205}).
\emph{Status: Complete.}
\end{proof}

\begin{proposition}[Changing type at a conjugate pair]\label{prop:P5-local}
\lean{Formalism.DivisionValues.ConjugateWeightGram.local_idW}
\uses{lem:P5-dual,lem:P5-ratio}
For fixed $y$ put
\[
 U_y(t)=v_y(G_{t,y})
       +\sum_j\{i_y(M_j,B_{jt})+i_{\bar y}(M_j,B_{jt})+j_y(H_{jt},T_{jt})\}.
\]
Then $U_y(t_y)-U_y(t_0)=i_y(A_{t_0},A_{t_y})+i_{\bar y}(A_{t_0},A_{t_y})$.
The same formula holds in the general determinant comparison when at
$v_0$ the type $t_y$ has scalar $\kappa_{v_0}\rho_r$ with
$\rho_r=\sigma_0(\rho)>0$, $v_y(\rho)=0$, and the base multiplier is one.
\end{proposition}
\begin{proof}
Use subscripts $0,1$ for $t_0,t_y$. Choose semilocal column bases
$f_{jh}$ for $T_{jh}$ and write $g_{jh}=\det(f_{jh}^*H_{jh} f_{jh})$.
Let $C$ be the change of basis from $\iota_0e_0$ to $\iota_1e_1$ and
$D_j$ the change from $\theta_{j0}f_{j0}$ to $\theta_{j1}f_{j1}$.
With $N(z)=z^*z$ these matrices give the following identity in $K$:
\[
  \rho^D G_1\prod_j\bigl\{N(\det D_j)g_{j0}\bigr\}
          =N(\det C)G_0\prod_j g_{j1}.                         \qquad\text{(P5.2)}
\]
To see it, first work at $\sigma_0$. The image row Gram determinants in
the common metric are $d_{v_0}(\iota_h)\sigma_0(G_h)$, and the ratio of
these image determinants is $|\sigma_0(\det C)|^2$. For each column the
corresponding equation is
\[
 d_{v_0}(\theta_{j1})\sigma_0(g_{j1})=
 |\sigma_0(\det D_j)|^2d_{v_0}(\theta_{j0})\sigma_0(g_{j0}).
\]
Insert these equations in the row equation, and divide by
$\kappa_{v_0}^D\prod_j d_{v_0}(\theta_{j0})$. It is nonzero by positivity
and invertibility. Since $\sigma_0$ respects the involution on $K$,
the result is the image of (P5.2); injectivity lifts it back to $K$.
Only this embedding occurs here. In particular the argument never applies
a finite valuation to a complex number.
See \path{ConjugateWeightGram.local_idW}
(\path{ConjugateWeightLatticeOfGeneral.lean:208--265}).

Every factor in (P5.2) is nonzero. Take $v_y$; the multiplier contributes
$D v_y(\rho)=0$. The valuation of $N(\det D_j)$ is the sum of the two
valuations of $\det D_j$ at $y,\bar y$. Since the same matrices represent
the two local changes of basis, these valuations are the two changes of
indices with reference lattice $M_j$. The term
$v_y(g_{j0})-v_y(g_{j1})$ is the \emph{increase} of the Hermitian
dual index by Lemma~\ref{lem:P5-dual}. Thus the difference of $U$ is
$v_y(\det C)+v_{\bar y}(\det C)$, as claimed. If $y=\bar y$ the two ordinary
indices coincide and both occurrences remain. If the type is unchanged both
sides vanish even when base Gram and column terms themselves are nonzero.
See \path{ConjugateWeightGram.local_idW}
(\path{ConjugateWeightLatticeOfGeneral.lean:266--317}).
\emph{Status: Complete.}
\end{proof}

